% =============================================================================
% ch05_theory_certificates.tex — Certificates: When the Boundary Can Be Crossed Safely
% Part of the K-Bound long-form monograph. NO preamble here; \input by main.tex.
% Faithful to docs/research/kbound/kbound.tex §5.3–§5.5 and the anytime appendix.
% =============================================================================

\chapter{Certificates: When the Boundary Can Be Crossed Safely}
\label{ch:theory2}

Chapter~\ref{ch:theory1} ended on deliberately negative ground. Theorem~\ref{thm:imp}
exhibited two target worlds that induce \emph{exactly} the same label-free evidence yet
assign \emph{opposite} signs to the adaptation benefit
$\Delta = R_T(f_0) - R_T(f_a)$, and its quantitative strengthening showed that the best
achievable committal error against such a pair equals $1-\TV$ of the evidence laws ---
which is $1$ in the witness. Theorem~\ref{thm:gate} then priced the failure: a gate's
regret over the oracle is precisely the benefit magnitude $|\Delta(Z)|$ it forfeits on
each sign mistake, and on evidence compatible with both signs no estimator can beat the
floor $\E[|\Delta(Z)|]/2$, attained by abstaining. The unknowable regime is not a
deficiency of any particular detector or drift statistic; it is a property of the
information available.

This chapter is the constructive half of the theory. It asks the question the
impossibility result leaves open: \emph{when the evidence does separate the worlds, how
should a system cross the boundary --- and what guarantee can it carry across?} The
answer is a \emph{certificate}: a label-free benefit estimate $\widehat\Delta(Z)$
equipped with a finite-sample confidence radius $\varepsilon(Z)$, together with the
three-way rule that commits only when the certified interval clears zero. We develop
the batch certificate and its error guarantee (Theorem~\ref{thm:cert}), discharge the
radius two ways --- by empirical-Bernstein concentration and by split-conformal
calibration --- and derive both the mandatory-abstention consequence on matched
evidence and the sample complexity of certification. We then build the streaming
complement: an anytime-valid certificate by testing-by-betting
(Theorem~\ref{thm:evalue}), valid at every stopping time simultaneously. The second
half of the chapter maps the regimes where the boundary genuinely opens: the
covariate-shift regime where the benefit sign becomes fully identifiable
(Theorem~\ref{thm:pos}), and the disagreement-region characterization that reduces
$\sign\Delta$ to an ordinal accuracy comparison on an observable set for binary,
multiclass, and regression losses
(Theorems~\ref{thm:disagree}--\ref{thm:disagree-reg}). The chapter closes with the one
piece the program leaves open --- the label-free \emph{bracketing} of that ordinal
comparison (Conjecture~\ref{conj:gen}) --- stated precisely, with an honest account of
why current techniques do not resolve it and which partial results are proved.

\section{From Impossibility to Conditional Possibility}
\label{sec:t2-idea}

The structure of the impossibility theorem dictates the structure of its escape. The
witness of Theorem~\ref{thm:imp} succeeds because the evidence law $\Law(Z)$ carries no
information about $\sign\Delta$: both worlds produce it identically. A certificate can
therefore exist only \emph{conditionally} --- under an explicit assumption that the
evidence at hand is not of the witness kind, i.e.\ that it constrains the benefit sign.
The paper names this property at the definition stage (Chapter~\ref{ch:formulation})
and invokes it as the standing assumption of the certificate result. We restate it
here, in the paper's own words, as the assumption the entire chapter trades under.

\begin{assumption}[Observable risk alignment, certificate form]
\label{sec:t2-ass-align}
The evidence is \emph{risk-aligned} in the sense of Chapter~\ref{ch:formulation}:
$Z$ identifies or bounds $\sign\Delta(z)$ (it need not reveal the label, only whether
adaptation helps). Operationally, a label-free estimator $\widehat\Delta(z)$ admits a
radius $\varepsilon(z)$ with
\[
\Pr\bigl[\,|\widehat\Delta(z)-\Delta(z)|\le \varepsilon(z)\,\bigr]\;\ge\;1-\alpha .
\]
\end{assumption}

The first sentence is the paper's definition of risk alignment, verbatim; the second is
the exact premise of the certificate theorem below. The two are the qualitative and
quantitative faces of the same requirement: a radius $\varepsilon(z)$ with valid
$1-\alpha$ coverage is precisely a device that \emph{bounds} $\sign\Delta(z)$ whenever
the interval $[\widehat\Delta-\varepsilon,\widehat\Delta+\varepsilon]$ excludes zero.

Given the assumption, the decision rule is forced by the geometry of the interval.
Estimate the benefit, attach the radius, and act on the \emph{certified} part of the
interval only:
\[
\textsc{adapt}\ \text{if}\ \widehat\Delta-\varepsilon>0,\qquad
\textsc{freeze}\ \text{if}\ \widehat\Delta+\varepsilon<0,\qquad
\textsc{abstain}\ \text{otherwise.}
\]
Adapting requires the entire interval to sit above zero; freezing requires it to sit
below; an interval that straddles zero is exactly the evidence state in which both
signs remain consistent with what has been observed, and the rule declines to commit.
The three outputs of the rule are the three regimes of the knowability map: knowably
helpful, knowably harmful, unknowable.

Before stating the theorem we are explicit about what its guarantee can and cannot
mean, because the certificate is easy to over-read.

\begin{remark}[What the guarantee means --- and what it does not]
\label{sec:t2-rem-meaning}
First, the guarantee is \emph{conditional} on
Assumption~\ref{sec:t2-ass-align}. Theorem~\ref{thm:imp} shows that the assumption
itself cannot be verified from $Z$ alone: on matched-evidence worlds no label-free
procedure can detect that the radius has lost its meaning. The certificate does not
contradict the impossibility theorem --- it relocates the burden into the validity of
$\varepsilon$, which in practice is purchased with a labeled validation or calibration
sample assumed exchangeable with deployment (the ``validation labels allowed'' sense of
label-free fixed in Chapter~\ref{ch:formulation}; no \emph{target} labels are ever
used). The paper is explicit that this residual assumption ``is not proved in full
generality,'' and we preserve that honesty here. Second, the guarantee bounds the
probability of the \emph{joint} event $\{\text{commit} \wedge \text{wrong sign}\}$; it
is not a posterior statement of the form ``given that we adapted, the chance of harm is
at most $\alpha$.'' Third, the certificate controls \emph{sign} errors, not magnitudes:
by Theorem~\ref{thm:gate} the regret of a wrong commitment is $|\Delta|$, so the
certificate's value concentrates exactly where $|\Delta|$ is large --- which is also
where the interval clears zero soonest. Fourth, abstention is not free; its price is
coverage, the benefit forgone when $\Delta>0$ but the interval straddles zero. The
level $\alpha$ is the dial that trades these two costs, and we will see the trade
measured on real tasks (Figure~\ref{fig:t2-alpha}).
\end{remark}

\section{The Finite-Sample Certificate}
\label{sec:t2-cert}

We now state the certificate theorem exactly as the paper does, and then expand its
proof to full length.

\begin{theorem}[Certificate]\label{thm:cert}
Suppose a label-free estimator $\widehat\Delta(z)$ admits a radius $\varepsilon(z)$ with
$\Pr[\,|\widehat\Delta(z)-\Delta(z)|\le \varepsilon(z)\,]\ge 1-\alpha$. Define
\[
\textsc{adapt}\ \text{if}\ \widehat\Delta-\varepsilon>0,\quad
\textsc{freeze}\ \text{if}\ \widehat\Delta+\varepsilon<0,\quad
\textsc{abstain}\ \text{otherwise.}
\]
Then $\Pr[\textsc{adapt}\wedge \Delta\le 0]\le\alpha$ and
$\Pr[\textsc{freeze}\wedge \Delta\ge 0]\le\alpha$.
\end{theorem}

\begin{proof}
Write $G:=\{\,|\widehat\Delta-\Delta|\le\varepsilon\,\}$ for the coverage event, so
that $\Pr[G]\ge 1-\alpha$ by hypothesis, and let $G^c$ denote its complement, of
probability at most $\alpha$. The argument is a containment of error events in $G^c$.

\emph{False adapt.} Suppose the rule outputs \textsc{adapt} and the coverage event $G$
holds. The decision \textsc{adapt} means $\widehat\Delta-\varepsilon>0$; the event $G$
means $\Delta\ge\widehat\Delta-\varepsilon$. Chaining the two inequalities,
\[
\Delta\;\ge\;\widehat\Delta-\varepsilon\;>\;0 ,
\]
so on $G$ an \textsc{adapt} decision is necessarily correct: $\Delta>0$. Contrapositively,
the false-adapt event $\{\textsc{adapt}\wedge\Delta\le0\}$ is disjoint from $G$, i.e.
\[
\{\textsc{adapt}\wedge\Delta\le0\}\;\subseteq\;G^c ,
\qquad\text{hence}\qquad
\Pr[\textsc{adapt}\wedge\Delta\le0]\;\le\;\Pr[G^c]\;\le\;\alpha .
\]

\emph{False freeze.} The argument is symmetric. On $G$ we also have
$\Delta\le\widehat\Delta+\varepsilon$, and the decision \textsc{freeze} means
$\widehat\Delta+\varepsilon<0$; chaining gives $\Delta<0$, so on $G$ a \textsc{freeze}
decision is necessarily correct. Hence
$\{\textsc{freeze}\wedge\Delta\ge0\}\subseteq G^c$ and its probability is at most
$\alpha$.

\emph{Abstention.} The third branch incurs no error event by construction: abstaining
asserts nothing about $\sign\Delta$, so there is nothing to be wrong about. Its cost is
accounted separately, as coverage. \qedhere
\end{proof}

\begin{remark}[One bad event covers both errors]
\label{sec:t2-rem-union}
Because both error events were shown to be subsets of the \emph{same} bad event $G^c$,
the proof in fact yields the slightly stronger statement
$\Pr[(\textsc{adapt}\wedge\Delta\le0)\,\cup\,(\textsc{freeze}\wedge\Delta\ge0)]\le\alpha$
when the radius is two-sided as assumed: there is no need to spend $\alpha$ twice. The
paper states the two one-sided bounds, which is the form the deployment guarantees
quote, and we keep that convention.
\end{remark}

The theorem is short because it is an exercise in honest bookkeeping: every gram of
statistical content lives in the premise, the validity of $(\widehat\Delta,\varepsilon)$.
The rest of this section discharges that premise in the two ways the system actually
uses --- a concentration-of-measure radius when $\widehat\Delta$ is a sample mean of
bounded paired benefits, and a split-conformal radius when $\widehat\Delta$ is an
arbitrary learned predictor of the benefit.

\subsection{Discharging the Radius I: Empirical-Bernstein Concentration}
\label{sec:t2-eb}

The paper's first instantiation applies when $\Delta$ is a mean of bounded
\emph{paired benefits}
\[
X_i \;=\; \ell(f_0(x_i),y_i)-\ell(f_a(x_i),y_i)\;\in\;[a,b]
\]
over a validated stress bank --- a calibration collection of conditions with held-out
labels (validation labels, never target labels) on which both the frozen and the
candidate model are evaluated. Then $\Delta=\E[X_i]$, the estimator
$\widehat\Delta=\tfrac1n\sum_iX_i$ is a bounded sample mean, and a lower confidence
bound on a bounded mean is a classical object. Two inequalities are relevant: the
distribution-free Hoeffding bound, and the variance-adaptive empirical-Bernstein bound
that the implementation uses.

\begin{lemma}[Hoeffding lower confidence bound \cite{hoeffding1963probability}]
\label{sec:t2-lem-hoeffding}
Let $X_1,\dots,X_n$ be independent with values in $[a,b]$, mean $\mu$, sample mean
$\widehat\mu$, and range $R=b-a$. For any $\alpha\in(0,1)$, with probability at least
$1-\alpha$,
\[
\mu\;\ge\;\widehat\mu - R\,\sqrt{\frac{\log(1/\alpha)}{2n}} .
\]
\end{lemma}

\begin{proof}
By Hoeffding's lemma, a bounded centered variable satisfies
$\E[e^{s(X_i-\mu)}]\le e^{s^2R^2/8}$ for all $s\in\R$. The Chernoff method applied to
$-\sum_i(X_i-\mu)$ gives
$\Pr[\widehat\mu-\mu\le -t]\le\exp(-2nt^2/R^2)$ for every $t>0$; setting the right side
equal to $\alpha$ and solving for $t$ yields
$t=R\sqrt{\log(1/\alpha)/(2n)}$, which is the displayed bound. \qedhere
\end{proof}

\begin{lemma}[Empirical-Bernstein lower confidence bound \cite{maurer2009empirical}]
\label{sec:t2-lem-eb}
Let $X_1,\dots,X_n$ ($n\ge2$) be i.i.d.\ with values in an interval of length $R=b-a$,
mean $\mu$, sample mean $\widehat\mu$, and unbiased sample variance $\widehat V$. For
any $\alpha\in(0,1)$, with probability at least $1-\alpha$,
\[
\mu\;\ge\;\widehat\mu
\;-\;\sqrt{\frac{2\widehat V\,\log(2/\alpha)}{n}}
\;-\;\frac{7R\,\log(2/\alpha)}{3(n-1)} .
\]
\end{lemma}

\begin{proof}[Proof sketch]
The result is Maurer and Pontil's empirical Bernstein bound, stated there for variables
in $[0,1]$; applying it to the rescaled variables $(X_i-a)/R$ and multiplying through
by $R$ gives the form above (the mean and radius scale by $R$, the variance by $R^2$).
The argument has two halves, each run at level $\alpha/2$ and combined by a union
bound. First, Bernstein's inequality with the \emph{true} variance $\sigma^2$ gives,
with probability at least $1-\alpha/2$,
$\mu\ge\widehat\mu-\sqrt{2\sigma^2\log(2/\alpha)/n}-R\log(2/\alpha)/(3n)$. Second, the
sample standard deviation concentrates around the true one: a self-bounding
concentration argument (the entropy method applied to the variance statistic) gives
$\sigma\le\sqrt{\widehat V}+R\sqrt{2\log(2/\alpha)/(n-1)}$ with probability at least
$1-\alpha/2$. Substituting the second display into the first and collecting the
$1/n$- and $1/(n-1)$-order terms produces the stated constant $7/3$; the bookkeeping is
exactly that of \cite{maurer2009empirical}. \qedhere
\end{proof}

Applying Lemma~\ref{sec:t2-lem-eb} to the paired benefits and to their negatives gives
a two-sided radius
\[
\varepsilon_{\mathrm{EB}}
\;=\;\sqrt{\frac{2\widehat V\,\log(2/\alpha)}{n}}
\;+\;\frac{7(b-a)\,\log(2/\alpha)}{3(n-1)} ,
\]
which satisfies the premise of Theorem~\ref{thm:cert} whenever the paired benefits are
independent or exchangeable draws from the deployment-relevant distribution. This is
the certificate implemented in the repository's
\texttt{switching\_certificate.py} and re-exported as
\texttt{kga/certificate.py::empirical\_bernstein} \cite{kbound2026repo}; the
implementation takes the range explicitly ($R=2$ for paired $|p-y|$ losses, since each
loss lies in $[0,1]$), and returns an \emph{infinite} radius when $n<2$ --- a
deliberate safety default under which the rule of Theorem~\ref{thm:cert} can only
abstain. The paper's analytic statements use the streamlined radius
$\varepsilon(n)=\sigma\sqrt{2\log(2/\alpha)/n}+3(b-a)\log(2/\alpha)/n$ with a variance
proxy $\sigma^2$ in place of $\widehat V$; the two differ only in constants and in the
empirical-versus-proxy variance, and we will use whichever form the quoted statement
uses.

\begin{remark}[Why empirical Bernstein, and why it is the right bound for this problem]
\label{sec:t2-rem-variance}
The Hoeffding radius scales with the \emph{range} $R$; the empirical-Bernstein radius
scales, in its leading term, with the \emph{standard deviation}. For paired benefits
this distinction is not a refinement but the essence of the problem. Under $0/1$ loss
the per-sample benefit $\delta=\ind[f_0\neq Y]-\ind[f_a\neq Y]$ takes values in
$\{-1,0,+1\}$ and vanishes identically off the disagreement region
$D=\{x:f_0(x)\neq f_a(x)\}$ --- the two losses cancel wherever the models agree, a fact
we exploit structurally in Section~\ref{sec:t2-disagree}. Consequently
$\E[\delta^2]\le P_T(D)$ (with equality in the binary case, where exactly one model is
correct on $D$), so the variance obeys $\mathrm{Var}(\delta)=\E[\delta^2]-\Delta^2\le P_T(D)$.
The empirical-Bernstein leading term is therefore on the order of
$\sqrt{2P_T(D)\log(2/\alpha)/n}$, while the Hoeffding radius is
$2\sqrt{\log(1/\alpha)/(2n)}$ regardless: when the candidate adaptation touches only a
small fraction of predictions --- the common case --- the variance-adaptive radius is
smaller by a factor of order $\sqrt{P_T(D)}$. All of the certificate's statistical
information is concentrated on the disagreement region, and the empirical-Bernstein
bound is the concentration inequality that knows it.
\end{remark}

We restate the paper's own honesty note before moving on. The empirical-Bernstein
radius is valid \emph{under independence or exchangeability of the calibrated benefit
sample with deployment}; in the paper's words, this assumption ``is the residual burden
and is not proved in full generality.'' Theorem~\ref{thm:cert} is proved; what remains
an assumption is that the world supplies data satisfying its premise. Chapter~\ref{ch:kga}
carries the same caveat into the algorithm's stated assumptions.

\subsection{Mandatory Abstention and the Cost of Certainty}
\label{sec:t2-forced}

Two corollaries calibrate expectations about the certificate, in opposite directions.
The first says that on genuinely unknowable evidence a valid certificate does not
merely \emph{tend} to abstain --- it \emph{must}. The second quantifies how much
calibration data the certificate needs before it can commit at all.

\begin{corollary}[Forced abstention under matched evidence]
\label{sec:t2-cor-abstain}
Fix evidence $z$ in the unknowable two-point regime of Theorem~\ref{thm:imp}
($\Law(Z\mid P_T^1)=\Law(Z\mid P_T^2)$, $\Delta^1<0<\Delta^2$). Any rule whose
false-adapt and false-freeze probabilities are each at most $\alpha$ (e.g.\ the
certificate of Theorem~\ref{thm:cert}) must abstain there with probability
$\Pr[\textsc{abstain}\mid z]\ge 1-2\alpha$.
\end{corollary}

\begin{proof}
A label-free rule observes only $Z$ (and independent internal randomness), so at the
fixed evidence value $z$ its action distribution is a single triple
$(q_a,q_f,1-q_a-q_f)$ --- the probabilities of \textsc{adapt}, \textsc{freeze}, and
\textsc{abstain} --- and, crucially, this triple is \emph{the same} in both worlds,
because the two evidence laws coincide. Now read the triple against each world in
turn. In world~$1$ the benefit is negative ($\Delta^1<0$), so every \textsc{adapt} the
rule emits at $z$ is a false adapt; the false-adapt constraint forces $q_a\le\alpha$.
In world~$2$ the benefit is positive ($\Delta^2>0$), so every \textsc{freeze} is a
false freeze; the false-freeze constraint forces $q_f\le\alpha$. The same $q_a$ and
$q_f$ must satisfy both constraints simultaneously, since the rule cannot tell the
worlds apart. Hence
$\Pr[\textsc{abstain}\mid z]=1-q_a-q_f\ge1-2\alpha$. \qedhere
\end{proof}

The paper's gloss deserves repeating: this makes abstention \emph{mandatory, not merely
conservative} --- on genuinely unknowable evidence a valid certificate cannot commit on
more than a $2\alpha$ fraction of decisions. The corollary is also the precise sense in
which the certificate and the impossibility theorem are two sides of one boundary: the
same $\alpha$ that licenses commitment where evidence separates \emph{caps} commitment
where it does not. Empirically, on the clean non-identifiability witness of
Chapter~\ref{ch:exp1} --- a constructed matched-evidence pair --- \KGA\ abstains on
$100\%$ of instances, while a rule forced to commit pays near-chance regret $0.51$;
Corollary~\ref{sec:t2-cor-abstain} is the theorem that says it had no valid
alternative.

\begin{corollary}[Sample complexity of certification]
\label{sec:t2-cor-n}
Let $\widehat\Delta$ be a mean of $n$ bounded paired benefits in $[a,b]$ with variance
proxy $\sigma^2$, and let
$\varepsilon(n)=\sigma\sqrt{2\log(2/\alpha)/n}+\tfrac{3(b-a)\log(2/\alpha)}{n}$
be the empirical-Bernstein radius of Theorem~\ref{thm:cert}. Then the certificate is
non-abstaining (commits \textsc{adapt} or \textsc{freeze}) once
\[
n\;\gtrsim\;\frac{2\sigma^2}{\Delta^2}\,\log\frac{2}{\alpha}\;+\;\frac{3(b-a)}{|\Delta|}\,\log\frac{2}{\alpha}
\;=\;O\!\Big(\tfrac{\sigma^2}{\Delta^2}\log\tfrac1\alpha\Big),
\]
so the unlabeled paired-benefit sample size needed to certify scales inversely with the
squared benefit margin and logarithmically with the confidence level.
\end{corollary}

\begin{proof}
The rule of Theorem~\ref{thm:cert} commits precisely when the interval clears zero,
i.e.\ when $\varepsilon(n)<|\widehat\Delta|$; since $\widehat\Delta\to\Delta$, the
asymptotically operative condition is $\varepsilon(n)<|\Delta|$. The radius is a sum of
two decreasing terms, so (up to the constant absorbed in $\gtrsim$) it suffices to make
each term smaller than $|\Delta|$. For the variance term,
\[
\sigma\sqrt{\frac{2\log(2/\alpha)}{n}}\;<\;|\Delta|
\quad\Longleftrightarrow\quad
n\;>\;\frac{2\sigma^2}{\Delta^2}\,\log\frac{2}{\alpha},
\]
which is the dominant requirement when $\sigma\gg|\Delta|$. Retaining the Bernstein
$1/n$ term adds the second requirement
$3(b-a)\log(2/\alpha)/n<|\Delta|$, i.e.\ $n>3(b-a)|\Delta|^{-1}\log(2/\alpha)$, the
stated summand. The asymptotic form follows because the first term dominates whenever
$\sigma^2/|\Delta|\gtrsim(b-a)$, which holds in the small-margin regime where the bound
bites. \qedhere
\end{proof}

The corollary answers what the paper calls the reviewer question --- \emph{when does
the certificate actually fire?} Certification is cheap where it matters most: large
benefit margins (precisely the decisions whose regret content is large, by
Theorem~\ref{thm:gate}) are certified from few samples, while margins near zero demand
$n\sim\sigma^2/\Delta^2$ and below that threshold the rule abstains --- which, by
Theorem~\ref{thm:gate} again, is exactly the band where committal regret is bounded by
the small $|\Delta|$ itself. The prediction is directly visible in the batch-size
ablation of Chapter~\ref{ch:exp1}: growing the test batch from $16$ to $256$ samples
shrinks the radius from $0.44$ to $0.18$ and raises coverage from $2\%$ to $26\%$,
the empirical shadow of the $n^{-1/2}$ term.

\subsection{Discharging the Radius II: The Split-Conformal Variant}
\label{sec:t2-conformal}

The empirical-Bernstein route requires $\widehat\Delta$ to be a sample mean of bounded
benefits. In the deployed system the benefit estimator is richer: \KGA\ trains a
gradient-boosted regressor mapping the evidence vector $Z$ to a predicted benefit,
leave-one-task-out across a bank of calibration tasks (Chapter~\ref{ch:kga}). For an
arbitrary learned $\widehat\Delta$, concentration inequalities no longer apply --- but
a conformal wrapper restores a finite-sample radius under a single exchangeability
assumption, with no further conditions on the estimator
\cite{vovk2005algorithmic,angelopoulos2023conformal}.

\begin{proposition}[Split-conformal radius]
\label{sec:t2-prop-conformal}
Let $(\widehat\Delta_1,\Delta_1),\dots,(\widehat\Delta_n,\Delta_n)$ be calibration
pairs and $(\widehat\Delta_{n+1},\Delta_{n+1})$ a deployment pair such that the
residuals $r_i=|\widehat\Delta_i-\Delta_i|$, $i=1,\dots,n+1$, are exchangeable. Let
$\varepsilon$ be the $\lceil(1-\alpha)(n+1)\rceil$-th smallest of $r_1,\dots,r_n$
(setting $\varepsilon=+\infty$ if $\lceil(1-\alpha)(n+1)\rceil>n$). Then
\[
\Pr\bigl[\,|\widehat\Delta_{n+1}-\Delta_{n+1}|\le\varepsilon\,\bigr]\;\ge\;1-\alpha ,
\]
so the rule of Theorem~\ref{thm:cert}, run with $(\widehat\Delta_{n+1},\varepsilon)$,
has false-adapt and false-freeze probabilities at most $\alpha$.
\end{proposition}

\begin{proof}
Set $k=\lceil(1-\alpha)(n+1)\rceil$ and let $r_{(1)}\le\dots\le r_{(n)}$ denote the
ordered calibration residuals, so $\varepsilon=r_{(k)}$ when $k\le n$. By
exchangeability of $(r_1,\dots,r_{n+1})$, the rank of $r_{n+1}$ within the full set of
$n+1$ residuals is stochastically dominated by the uniform distribution on
$\{1,\dots,n+1\}$ (uniform under almost-sure distinctness; ties only push the rank
down, which helps). The event $r_{n+1}\le r_{(k)}$ contains the event that the rank of
$r_{n+1}$ is at most $k$, whence
\[
\Pr\bigl[r_{n+1}\le r_{(k)}\bigr]\;\ge\;\frac{k}{n+1}
\;=\;\frac{\lceil(1-\alpha)(n+1)\rceil}{n+1}\;\ge\;1-\alpha .
\]
If $k>n$ the radius is infinite and the coverage statement is trivial. The premise of
Theorem~\ref{thm:cert} is thus satisfied at level $\alpha$ for the deployment pair, and
its conclusion follows. \qedhere
\end{proof}

The guarantee is marginal over the calibration draw --- the standard split-conformal
semantics \cite{angelopoulos2023conformal} --- which matches exactly the marginal
reading of Theorem~\ref{thm:cert} discussed in Remark~\ref{sec:t2-rem-meaning}. Two
implementation notes keep the record straight. First, the experiments use the empirical
$(1-\alpha)$-quantile of leave-one-task-out residuals
(\texttt{kga/certificate.py::conformal\_split}, mirroring the knowability and
mixed-regime experiment scripts \cite{kbound2026repo}); this differs from the
$\lceil(1-\alpha)(n+1)\rceil$ order statistic by an index shift of order $1/n$, the
usual practical simplification. Second, the exchangeability here is \emph{across
tasks}: the calibration bank of tasks and the deployment task are assumed exchangeable.
The paper lists this as a standing limitation, and Chapter~\ref{ch:exp1} probes it
empirically rather than assuming it away. Altogether the repository ships four
interchangeable radius constructions behind one interface --- empirical-Bernstein,
Hoeffding, split-conformal, and the anytime e-value process of the next section ---
all pure \texttt{numpy} and deterministic \cite{kbound2026repo}.

\begin{figure}[t]\centering
\includegraphics[width=0.78\textwidth]{fig_certificate.png}
\caption{The batch certificate on the 123-task anomaly-routing suite
(Chapter~\ref{ch:exp1}). Each point is one task: label-free estimated benefit
$\widehat\Delta(Z)$ (horizontal) against the oracle true benefit (vertical), with the
split-conformal band $\pm0.049$ at $\alpha=0.1$ shaded around the ideal diagonal.
Green points are certified \textsc{adapt}, red \textsc{freeze}, grey \textsc{abstain}:
commitments concentrate where the certified interval clears zero, and abstentions
cluster in the near-zero-benefit band where Theorem~\ref{thm:gate} prices commitment
errors as cheap but Theorem~\ref{thm:cert} declines to make them.}
\label{fig:t2-certificate}
\end{figure}

\begin{figure}[t]\centering
\includegraphics[width=0.55\textwidth]{fig_false_adapt.png}
\caption{Suppression of harmful adaptation on the 123-task suite (produced by the
knowability experiment of Chapter~\ref{ch:exp1}; same data as
Figure~\ref{fig:t2-certificate}). Each bar is the fraction of adaptation decisions
whose true benefit was negative: adapting always inherits the suite's harmful base
rate $0.358$; a naive validation-threshold rule false-adapts on $0.180$ of its
commitments; the certificate trichotomy on $0.100$ --- the complement of its $0.90$
adapt precision, in line with its calibrated level $\alpha=0.1$.}
\label{fig:t2-false-adapt}
\end{figure}

Figures~\ref{fig:t2-certificate} and~\ref{fig:t2-false-adapt} show the certificate
doing on real tasks what the theorem promises on paper. On the 123-task suite the
conformal band has half-width $0.049$; tasks whose interval clears zero are committed,
and the adapt precision is $0.90$ there and $0.935$ in the controlled mixed regime ---
in the latter, an observed false-adapt rate of $0.065$ against the budget
$\alpha=0.1$. The bar chart makes the comparative point on the same 123 tasks: an
uncertified always-adapt policy inherits the suite's $35.8\%$ harmful base rate, a
naive validation threshold false-adapts on $18.0\%$ of its commitments, and the
trichotomy holds the line at $10.0\%$ --- the level it was asked to hold.

\section{The Anytime-Valid Certificate: Testing by Betting}
\label{sec:t2-evalue}\label{app:anytime}

Theorem~\ref{thm:cert} fixes the sample size $n$ in advance. Streaming test-time
adaptation does not: evidence arrives one batch at a time, the system looks at the
accumulating benefit sample after every batch, and it wants to commit \emph{the moment}
the evidence suffices --- without scheduling $n$ beforehand and without paying a
multiple-testing tax for the repeated looks. Classical fixed-$n$ confidence intervals
are unsound under this regime: stopping the first time a fixed-$n$ interval clears zero
inflates the error probability well above $\alpha$, because the analyst's stopping rule
correlates with the noise. What is needed is a certificate whose guarantee holds
\emph{simultaneously over all stopping times}. The paper supplies one in its appendix,
via the modern theory of safe anytime-valid inference
\cite{ramdas2023gametheoretic,howard2021timeuniform,grunwald2024safe}, in the
testing-by-betting form popularized by Shafer \cite{shafer2019game} with the betting
strategies of Waudby-Smith and Ramdas \cite{waudbysmith2024betting}.

\subsection{Wealth, Supermartingales, and Ville's Inequality}
\label{sec:t2-betting}

Testing by betting replaces the $p$-value with a gambler's wealth. To test the null
hypothesis $H_0:\Delta\le0$ (``adapting does not help''), imagine a skeptic who starts
with capital $1$ and, before each benefit observation $X_t$, stakes a fraction of
current wealth on the proposition that the benefits are drifting positive. If the bet
at round $t$ is $\lambda_t\ge0$, the wealth multiplies by $(1+\lambda_tX_t)$: it grows
when $X_t>0$, shrinks when $X_t<0$. Under the null, every such bet is fair or
unfavorable --- the conditional expected multiplier is
$1+\lambda_t\E[X_t\mid\mathcal F_{t-1}]\le1$ --- so no betting system can be expected
to multiply the skeptic's capital. Large realized wealth is therefore quantitative
evidence against the null: wealth $E^+_t\ge1/\alpha$ is the event that a gambler
playing a fair-or-losing game has multiplied capital ten-fold (at $\alpha=0.1$), and
the probability that this \emph{ever} happens is at most $\alpha$. The wealth process
is an \emph{e-process} --- a nonnegative process whose value at any stopping time has
expectation at most one under the null \cite{grunwald2024safe,ramdas2023gametheoretic}
--- and the ``ever'' quantifier is exactly what frees the certificate from fixed-$n$
scheduling. The inequality that delivers the quantifier is Ville's
\cite{ville1939etude}.

\begin{lemma}[Ville's inequality \cite{ville1939etude}]
\label{sec:t2-lem-ville}
Let $(M_t)_{t\ge0}$ be a nonnegative supermartingale with respect to a filtration
$(\mathcal F_t)$, with $\E[M_0]\le1$. Then for any $u\ge1$,
\[
\Pr\Bigl[\,\exists\,t\ge0:\;M_t\ge u\,\Bigr]\;\le\;\frac{1}{u} .
\]
\end{lemma}

\begin{proof}
Let $\tau:=\inf\{t\ge0:M_t\ge u\}$, a stopping time. Fix a horizon $T$; then
$\tau\wedge T$ is a bounded stopping time, and the optional stopping theorem for
nonnegative supermartingales gives $\E[M_{\tau\wedge T}]\le\E[M_0]\le1$. On the event
$\{\tau\le T\}$ we have $M_{\tau\wedge T}=M_\tau\ge u$, while on its complement
$M_{\tau\wedge T}\ge0$; hence
\[
1\;\ge\;\E[M_{\tau\wedge T}]\;\ge\;u\,\Pr[\tau\le T] ,
\qquad\text{so}\qquad \Pr[\tau\le T]\;\le\;\frac1u .
\]
Letting $T\to\infty$ and using continuity from below,
$\Pr[\tau<\infty]\le1/u$, which is the claim. \qedhere
\end{proof}

Markov's inequality bounds $\Pr[M_t\ge u]$ at one preselected $t$; Ville's inequality
bounds the supremum over the whole trajectory at no extra cost. That uniformity over
time is the entire content of ``anytime-valid'': the guarantee survives optional
stopping, optional continuation, and continuous monitoring
\cite{howard2021timeuniform}.

\subsection{The Anytime Certificate}
\label{sec:t2-evalue-thm}

We state the theorem as the paper's appendix does. Two wealth processes run in
parallel: $E^+_t$ bets that the mean benefit is positive (its growth rejects
$H_0:\Delta\le0$ and certifies \textsc{adapt}), and $E^-_t$ bets that it is negative
(its growth rejects $H_0':\Delta\ge0$ and certifies \textsc{freeze}); while neither has
grown past $1/\alpha$, the system abstains --- which in the streaming reading means
``keep sampling.''

\begin{theorem}[Anytime-valid adapt/freeze/abstain certificate]
\label{thm:evalue}
Let paired benefits $X_i=\ell(f_0(x_i),y_i)-\ell(f_a(x_i),y_i)\in[a,b]$ ($a<0<b$) be
adapted to $(\mathcal F_t)$ with $\E[X_t\mid\mathcal F_{t-1}]=\Delta$. With predictable
bets $\lambda_t\in[0,c/(-a)]$, $\nu_t\in[0,c/b]$, $c\in(0,1)$, set
$E_t^{+}=\prod_{i\le t}(1+\lambda_i X_i)$ and $E_t^{-}=\prod_{i\le t}(1-\nu_i X_i)$.
Decide \textsc{adapt} the first time $E_t^{+}\ge1/\alpha$, \textsc{freeze} the first time
$E_t^{-}\ge1/\alpha$, else \textsc{abstain}. Then simultaneously over all $t$,
$\Pr[\exists t: E_t^{+}\ge1/\alpha \wedge \Delta\le0]\le\alpha$ and
$\Pr[\exists t: E_t^{-}\ge1/\alpha \wedge \Delta\ge0]\le\alpha$.
\end{theorem}

\begin{proof}
We prove the \textsc{adapt} direction in full; the \textsc{freeze} direction is its
mirror image.

\emph{Step 1: the wealth stays positive.} Each factor of $E^+_t$ is
$1+\lambda_iX_i$ with $\lambda_i\ge0$ and $X_i\ge a$, so the worst case is $X_i=a<0$,
where the factor is at least $1+\lambda_ia\ge1+\tfrac{c}{-a}\,a=1-c>0$ by the bet cap.
A product of strictly positive factors is strictly positive, so $E^+_t>0$ for all $t$;
in particular $E^+_t$ is a legitimate candidate for Lemma~\ref{sec:t2-lem-ville}.

\emph{Step 2: under the null, the wealth is a supermartingale.} Fix any distribution in
$H_0$, i.e.\ with $\Delta\le0$. Because $\lambda_t$ is \emph{predictable} --- it is
$\mathcal F_{t-1}$-measurable, a function of the past only --- it factors out of the
conditional expectation:
\[
\E\bigl[E^+_t\,\big|\,\mathcal F_{t-1}\bigr]
=E^+_{t-1}\,\E\bigl[1+\lambda_tX_t\,\big|\,\mathcal F_{t-1}\bigr]
=E^+_{t-1}\,\bigl(1+\lambda_t\,\E[X_t\mid\mathcal F_{t-1}]\bigr)
=E^+_{t-1}\,(1+\lambda_t\Delta)
\;\le\;E^+_{t-1},
\]
since $\lambda_t\ge0$ and $\Delta\le0$ make the multiplier at most one. With
$E^+_0=1$, the process $(E^+_t)$ is a nonnegative supermartingale with unit initial
value under every null distribution.

\emph{Step 3: Ville closes the bound.} Lemma~\ref{sec:t2-lem-ville} with $u=1/\alpha$
gives, under every $\Delta\le0$,
\[
\Pr\bigl[\exists t:\,E^+_t\ge1/\alpha\bigr]\;\le\;\alpha .
\]
The false-adapt event $\{\exists t:E^+_t\ge1/\alpha\wedge\Delta\le0\}$ is exactly this
crossing event evaluated under a null distribution, so its probability is at most
$\alpha$ --- uniformly over all stopping times, since the bound already quantifies over
the whole trajectory.

\emph{The freeze direction.} The process $E^-_t=\prod_{i\le t}(1-\nu_iX_i)$ has factors
$1-\nu_iX_i$ with $\nu_i\in[0,c/b]$; the worst case is now $X_i=b>0$, where the factor
is at least $1-\nu_ib\ge1-c>0$, so $E^-_t>0$. Under $H_0':\Delta\ge0$, predictability
of $\nu_t$ gives
$\E[E^-_t\mid\mathcal F_{t-1}]=E^-_{t-1}(1-\nu_t\Delta)\le E^-_{t-1}$, and Ville's
inequality bounds the crossing probability by $\alpha$ exactly as before. \qedhere
\end{proof}

\begin{remark}[Validity is free; power is bought by the bet]
\label{sec:t2-rem-bets}
Nothing in the proof used the \emph{values} of the bets beyond nonnegativity,
predictability, and the cap --- any such betting sequence yields a valid certificate.
The specific bets govern power and stopping time. The implementation uses a truncated
growth-optimal (aGRAPA-style) plug-in
\[
\lambda_t=\mathrm{clip}\Bigl(\widehat\mu_{t-1}\,/\,\widehat\sigma^2_{t-1},\;0,\;c/(-a)\Bigr),
\]
the running-mean-over-running-second-moment approximation to the log-optimal bet
$\lambda^\star=\Delta/\E[X^2]$ \cite{waudbysmith2024betting}, with a small prior
stabilizing the early variance estimate; the symmetric rule drives $\nu_t$. This
separation --- validity from the supermartingale property, power from the betting
strategy --- is the signature design freedom of game-theoretic statistics
\cite{shafer2019game,ramdas2023gametheoretic}. The anytime certificate
\emph{complements} rather than replaces the batch certificate: at a pre-fixed $n$ the
empirical-Bernstein interval is tighter, while the e-process is valid at any stopping
time; the paper states exactly this trade.
\end{remark}

\subsection{The Validator: Measured Anytime Error and Power}
\label{sec:t2-evalue-validation}

Theorem~\ref{thm:evalue} ships with a machine-checked numerical validator,
\texttt{experiments/kbound/theory\_validation/val\_thm3\_evalue.py}
\cite{kbound2026repo}, which runs the betting certificate over $20{,}000$ independent
streams per configuration (horizon $2000$, $\alpha=0.1$, bet cap $c=0.5$, seed $0$) and
three bounded benefit distributions chosen to stress different shapes within
$[a,b]=[-1,1]$: a two-point distribution massing on the endpoints (the worst case for
the positivity argument), a smooth Beta-shaped law, and a clipped Gaussian. It measures
the \emph{anytime} false-adapt rate --- the fraction of null streams in which $E^+_t$
\emph{ever} crosses $1/\alpha$ within the horizon --- together with detection power and
stopping times under alternatives. Table~\ref{sec:t2-tab-evalue} reproduces the
machine-readable results file.

\begin{table}[t]\centering\small
\begin{tabular}{lcccccc}
\toprule
& \multicolumn{3}{c}{anytime false-adapt ($H_0:\Delta\le0$)}
& \multicolumn{3}{c}{power / mean stop ($H_1:\Delta>0$)}\\
\cmidrule(lr){2-4}\cmidrule(lr){5-7}
stream & $\Delta=-0.20$ & $\Delta=-0.05$ & $\Delta=0$
& $\Delta=0.05$ & $\Delta=0.10$ & $\Delta=0.40$\\
\midrule
two-point        & $0.00505$ & $0.0268$  & $0.0626$ & $0.375\,/\,810$ & $0.948\,/\,676$ & $1.000\,/\,39$\\
Beta             & $0.0000$  & $0.0038$  & $0.0584$ & $0.988\,/\,521$ & $1.000\,/\,127$ & $1.000\,/\,16$\\
clipped Gaussian & $0.0000$  & $0.00465$ & $0.0559$ & $0.976\,/\,588$ & $1.000\,/\,148$ & $1.000\,/\,17$\\
\bottomrule
\end{tabular}
\caption{Anytime-valid certificate validation, exactly as recorded in
\texttt{experiments/kbound/theory\_validation/val\_thm3\_evalue\_results.json}
($\alpha=0.1$, wealth threshold $1/\alpha=10$, horizon $2000$, $20{,}000$ Monte Carlo
runs per cell, seed $0$, bet cap $c=0.5$; power rounded to three decimals, mean
stopping times to integers). The worst-case anytime false-adapt rate across all nine
null settings is $0.0626\le\alpha=0.1$; the results file's verdict field records
\texttt{controlled:\,true}.}
\label{sec:t2-tab-evalue}
\end{table}

Three features of the table are worth reading off. First, the headline: the worst-case
anytime false-adapt rate over every null configuration is $\mathbf{0.0626}$, against
the budget $\alpha=0.1$ --- the guarantee of Theorem~\ref{thm:evalue} holds with room
to spare, and the worst case occurs exactly where theory says it must, at the boundary
$\Delta=0$ of the null (rates collapse to $0.00505$ and below by $\Delta=-0.20$). The
paper's appendix additionally records a run at $\alpha=0.05$ with anytime false-adapt
$0.032\le0.05$. The slack between $0.0626$ and $0.1$ is the familiar conservatism of
time-uniform bounds: the certificate pays for validity at \emph{every} stopping time,
and the plug-in bet is not the oracle log-optimal one. Second, the supermartingale
property itself is checked directly: at $\Delta=0$ the Monte Carlo mean of $E^+_t$
stays at or below $1$ (within three standard errors) at every recorded time
$t\in\{1,5,25,100,500,2000\}$ for all three streams --- the results file flags
\texttt{all\_leq\_one\_within\_3sem:\,true} in each case. Third, power behaves as a
sequential test should: it rises to $1$ as the margin grows, and the mean stopping time
falls steeply --- from roughly $810$ samples at $\Delta=0.05$ (two-point stream) to
$16$--$39$ samples at $\Delta=0.40$. Large benefits are certified almost immediately;
marginal ones take evidence; null ones, with probability at least $1-\alpha$, are never
falsely certified no matter how long the system watches.

\begin{remark}[Bounded non-stationarity]
\label{sec:t2-rem-drift}
The constant-mean premise $\E[X_t\mid\mathcal F_{t-1}]=\Delta$ can be relaxed. The
paper's theory-extensions section proves a $\beta$-bounded-drift variant: if the
conditional mean wanders within $\pm\beta$ of $\Delta$, dividing the wealth by
$\prod_{i\le t}(1+\lambda_i\beta)$ restores the supermartingale property under the
null, at the price of an inflated decision threshold
$(1/\alpha)\prod_{i\le t}(1+\lambda_i\beta)\le\exp(\beta\sum_{i\le t}\lambda_i)$.
Its validator (\texttt{val\_thm9prime\_drift.py} \cite{kbound2026repo}) reports that
under worst-case upward drift the corrected certificate holds the anytime false-adapt
rate at $0.05/0.04/0.01$ for $\beta=0.1/0.2/0.4$ --- all within $\alpha=0.1$ --- while
the uncorrected threshold's rate inflates to $0.69$ and beyond, confirming both that
the correction works and that it is needed.
\end{remark}

\begin{figure}[t]\centering
\includegraphics[width=0.6\textwidth]{fig_alpha_coverage.png}
\caption{The level $\alpha$ as the safety--coverage dial, measured on the
anomaly-routing benchmark of Chapter~\ref{ch:exp1} with the split-conformal batch
certificate: $\alpha=0.01$ yields a $0\%$ observed false-adapt rate at $11\%$ coverage;
$\alpha=0.1$ yields $0.065$ at $35\%$ coverage; $\alpha=0.2$ yields $0.072$ at $49\%$
coverage. The $\alpha=0.05$ point sits at $0.073$, slightly above its nominal budget
--- an honest finite-sample fluctuation on $369$ instances under approximate cross-task
exchangeability, and a reminder that the guarantee is marginal, not per-configuration.}
\label{fig:t2-alpha}
\end{figure}

Both certificates --- batch and anytime --- expose the same single dial, the level
$\alpha$, and Figure~\ref{fig:t2-alpha} shows what the dial buys on real tasks.
Tightening $\alpha$ to $0.01$ drives observed false adaptation to zero at the cost of
committing on only $11\%$ of decisions; loosening it to $0.2$ roughly quadruples
coverage while the observed false-adapt rate stays near $0.07$. The certificate makes
the safety--coverage trade explicit and tunable, rather than implicit in a mechanism's
hyperparameters.

\section{A Provably Knowable Regime: Covariate Shift}
\label{sec:t2-pos}

The impossibility theorem and the certificate bracket the problem from both sides, but
neither yet exhibits a natural class of shifts on which the label-free decision is
\emph{provably} available. The paper's fourth result identifies one: covariate shift.
The motivating picture is a new scanner --- image statistics change, disease prevalence
and the disease-given-image relationship do not. Under that premise, reweighting
labeled \emph{source} data by the density ratio of the two covariate laws reconstructs
every target risk, and with it the benefit and its sign.

\begin{theorem}[Covariate-shift identifiability]
\label{thm:pos}
Assume covariate shift $P_S(Y\mid X)=P_T(Y\mid X)$, and a density ratio
$r(x)=p_T(x)/p_S(x)$ that is bounded ($r\le B<\infty$) with $\E_S[r(X)^2]<\infty$.
Then $R_T(f)=\E_S[r(X)\,\ell(f(X),Y)]$, the importance-weighted estimator
$\widehat R_T(f)=\tfrac1n\sum_i r(x_i)\ell(f(x_i),y_i)$ is unbiased and consistent, and
hence $\Delta$ and $\sign\Delta$ are identifiable from labeled source plus unlabeled
target whenever $|\Delta|>0$.
\end{theorem}

\begin{proof}
\emph{Change of measure.} Write the target risk as an integral over the target joint
and factor the joint through the covariate-shift premise:
\[
R_T(f)=\E_{P_T}\bigl[\ell(f(X),Y)\bigr]
=\!\int\!\!\int \ell(f(x),y)\,p_T(y\mid x)\,p_T(x)\,dy\,dx
=\!\int\!\!\int \ell(f(x),y)\,p_S(y\mid x)\,\frac{p_T(x)}{p_S(x)}\,p_S(x)\,dy\,dx ,
\]
where the second equality substitutes $p_T(y\mid x)=p_S(y\mid x)$ (the premise) and
multiplies and divides by $p_S(x)$; boundedness of $r$ ensures the manipulation is
legitimate wherever $p_S$ has mass, and the target law places no mass where $p_S$
vanishes (else $r$ would be unbounded). The right-hand side is
$\E_S[r(X)\,\ell(f(X),Y)]$, the first claim.

\emph{Unbiasedness and consistency.} Unbiasedness of $\widehat R_T(f)$ is immediate
from the identity just proved, since each summand $r(x_i)\ell(f(x_i),y_i)$ has
source-expectation $R_T(f)$. For consistency, the summands are i.i.d.\ with finite
second moment: with $\ell$ bounded in $[0,1]$,
$\E_S[(r\,\ell)^2]\le\E_S[r^2]<\infty$, so the strong law of large numbers gives
$\widehat R_T(f)\to R_T(f)$ almost surely, with variance
$\mathrm{Var}(\widehat R_T(f))\le\E_S[r^2]/n=O(1/n)$.

\emph{Sign identifiability.} Apply the estimator to both models and difference:
$\widehat\Delta=\widehat R_T(f_0)-\widehat R_T(f_a)
=\tfrac1n\sum_ir(x_i)\bigl(\ell(f_0(x_i),y_i)-\ell(f_a(x_i),y_i)\bigr)$ is unbiased
and consistent for $\Delta$ by linearity. If $|\Delta|>0$, then almost surely
$\sign\widehat\Delta=\sign\Delta$ for all large $n$, so the sign --- and indeed the
value --- of the benefit is identified from labeled source data plus the unlabeled
target sample needed to estimate $r$. \qedhere
\end{proof}

\paragraph{Why this escapes the impossibility theorem.}
It is worth dwelling on \emph{where} the witness of Theorem~\ref{thm:imp} dies under
the covariate-shift premise, because the mechanism is the template for every positive
result in this theory. The witness constructed two worlds agreeing on $P_T(X)$ and
differing only in $P_T(Y\mid X)$ --- the one coordinate of the target joint that
label-free evidence cannot see. Covariate shift removes exactly that degree of freedom:
it pins $P_T(Y\mid X)$ to $P_S(Y\mid X)$, an object estimable from \emph{labeled
source} data. The accessible data then determine the entire target joint,
$p_T(x,y)=p_T(x)\,p_S(y\mid x)$, with $p_T(x)$ read off the unlabeled target stream and
$p_S(y\mid x)$ off the source. The evidence has become \emph{sufficient} for $\Delta$.
Equivalently: the two witness worlds have different conditionals, so at most one of
them can satisfy the covariate-shift premise with respect to any fixed source ---
inside the assumption class, matched-evidence pairs with opposite benefit simply do not
exist. Identifiability results of this shape are the classical core of the
covariate-shift literature \cite{shimodaira2000improving,sugiyama2007covariate,
quinonero2009dataset}; the delta here is only the target of inference, the \emph{sign
of a risk difference} rather than a risk.

\paragraph{The price: weight boundedness and effective sample size.}
The premise buys identifiability; the moment conditions buy a usable certificate. The
weighted paired benefits $r(x_i)\bigl(\ell(f_0)-\ell(f_a)\bigr)$ are bounded in
$[-B,B]$ when $r\le B$, so they can be fed directly to the empirical-Bernstein
certificate of Section~\ref{sec:t2-eb}, and the radius inherits the weight
distribution's second moment. Since $\E_S[r]=1$, the variance inflation relative to
unweighted estimation is governed by $\E_S[r^2]\ge1$, with equality only under no
shift; the corresponding \emph{effective sample size} is $n/\E_S[r^2]$ (empirically,
$(\sum_iw_i)^2/\sum_iw_i^2$), and the certificate radius scales like
$\sqrt{\E_S[r^2]/n}$. Heavy-tailed weights therefore do not silently corrupt the
decision --- they \emph{visibly inflate the radius}, and the certificate responds by
abstaining or freezing. This is not a hypothetical failure mode: in the regression
covariate-shift experiment of Chapter~\ref{ch:exp1}, a $4\sigma$ shift makes the
importance-weighted candidate \emph{hurt} on $79\%$ of instances through sheer weight
variance, and \KGA\ refuses it wholesale ($0$ adapt, $11$ freeze, $61$ abstain),
matching the oracle's mean squared error of $0.251$ while always-adapt degrades to
$0.338$. Identifiability in principle and certifiability at a given $n$ are different
currencies, and the certificate trades only in the second.

\paragraph{The caveat that keeps the theorem honest.}
The paper attaches a caveat tying the result back to Theorem~\ref{thm:imp}, and it is
load-bearing: the premise $P_S(Y\mid X)=P_T(Y\mid X)$ is \emph{not} testable from
unlabeled data; when it fails, the regime can re-enter the unknowable region. Covariate
shift is an assumption about precisely the coordinate that label-free evidence cannot
observe, so adopting it is an act of modeling, not of inference. The knowability
hierarchy of Table~\ref{sec:t2-tab-shift} situates this regime among its neighbors:
label shift is similarly identifiable under its own structural premises (stable
class-conditionals and an invertible confusion matrix, the BBSE conditions
\cite{lipton2018detecting}); concept shift --- a change in $P(Y\mid X)$ itself --- is
non-identifiable and lands back in Theorem~\ref{thm:imp}; and the mixed or unknown case
is exactly where the certify-or-abstain machinery of this chapter operates.

\begin{table}[t]\centering\small
\begin{tabular}{lll}
\toprule
Shift type & Label-free status & Mechanism / theorem \\
\midrule
Covariate ($P(Y\!\mid\!X)$ fixed) & identifiable & importance weighting (Thm~\ref{thm:pos}) \\
Label / prior ($P(Y)$) & identifiable$^\dagger$ & stable class-conditionals, invertible confusion ($^\dagger$BBSE-type) \\
Disagreement-local & ordinal-identifiable & sign $=$ accuracy gap on $D$ (Thms~\ref{thm:disagree}--\ref{thm:disagree-reg}) \\
Concept ($P(Y\!\mid\!X)$ changes) & non-identifiable & re-enters Thm~\ref{thm:imp} \\
Mixed / unknown & certify-or-abstain & commit where $Z$ separates, else abstain \\
\bottomrule
\end{tabular}
\caption{Knowability hierarchy of distribution shifts for label-free adapt/freeze
decisions, as drawn in the paper. The boundary is evidence-identifiability: some shifts
are knowable, some are not, and a valid certificate abstains across the boundary. The
label-shift row follows the black-box shift estimation conditions of
\cite{lipton2018detecting}.}
\label{sec:t2-tab-shift}
\end{table}

\section{The Disagreement Region: Where the Sign of Benefit Lives}
\label{sec:t2-disagree}

The covariate-shift result buys identifiability with a global assumption on the shift.
The fifth result of the theory takes the opposite route: no assumption on the shift at
all, but a sharper look at \emph{where} in input space the decision is actually
decided. The observation is elementary and, once made, structural. Wherever the frozen
and candidate models agree, their losses are identical and the per-sample benefit is
exactly zero --- those points contribute \emph{nothing} to $\Delta$, under any label
law whatsoever. The entire benefit, sign and magnitude alike, is carried by the
\emph{disagreement region}
\[
D \;=\; \{x : f_0(x)\neq f_a(x)\},
\]
and $D$ is \emph{observable}: membership requires evaluating two known functions, no
labels. Restricted to $D$, the question ``does adapting help?'' collapses to an
above-chance check --- is the candidate right more often than the incumbent on the
points where they differ? --- which is an \emph{ordinal} comparison, not a risk
estimate.

\begin{theorem}[Sign-of-difference identifiability on the disagreement region; binary case]
\label{thm:disagree}
Let $Y\in\{0,1\}$ with $0/1$ loss, and let $D=\{x: f_0(x)\neq f_a(x)\}$ be the
\emph{disagreement region} (observable, since $f_0,f_a$ are known maps of $X$). Then
\[
\Delta \;=\; P_T(D)\,\bigl(2\,a_a^D-1\bigr),\qquad a_a^D:=P_T\!\left(f_a(X)=Y \mid X\in D\right),
\]
so $\sign\Delta=\sign\!\left(a_a^D-\tfrac12\right)$. Consequently $\sign\Delta$ is
identifiable from label-free evidence \emph{if and only if} the evidence determines
whether $f_a$ exceeds chance accuracy on $D$; in particular, if a label-free
reliability signal brackets $a_a^D$ relative to $\tfrac12$, the sign is certified.
\end{theorem}

\begin{proof}
Let $\delta(x,y)=\ell(f_0(x),y)-\ell(f_a(x),y)=\ind[f_0(x)\neq y]-\ind[f_a(x)\neq y]$
denote the per-sample benefit, so that $\Delta=\E_{P_T}[\delta]$.

\emph{Off $D$, the benefit vanishes identically.} If $x\notin D$ then
$f_0(x)=f_a(x)$, so the two indicators coincide for every $y$ and $\delta(x,y)=0$
pointwise --- not merely in expectation, but for each realization of the label.

\emph{On $D$, the benefit is a fair-coin score for $f_a$.} If $x\in D$ the two
predictions differ, and since $Y$ is binary, exactly one of $f_0(x),f_a(x)$ equals
$y$. If $f_a(x)=y$ then $f_0(x)\neq y$, so $\delta=1-0=+1$; if $f_0(x)=y$ then
$\delta=0-1=-1$. Hence, conditionally on $D$, $\delta$ is a $\pm1$ variable equal to
$+1$ with probability $a_a^D$, and
\[
\E[\delta\mid D]
=P_T(f_a{=}Y\mid D)-P_T(f_0{=}Y\mid D)
=a_a^D-\bigl(1-a_a^D\bigr)
=2a_a^D-1 ,
\]
using that on $D$ the two correctness events are complementary,
$P_T(f_0{=}Y\mid D)=1-a_a^D$.

\emph{Total expectation.} Splitting $\Delta$ over the partition $\{D,D^c\}$ and using
$\delta\equiv0$ off $D$,
\[
\Delta=\E\bigl[\delta\,\ind\{X\in D\}\bigr]+\E\bigl[\delta\,\ind\{X\notin D\}\bigr]
=P_T(D)\,\E[\delta\mid D]
=P_T(D)\,(2a_a^D-1).
\]
Since $P_T(D)\ge0$, the sign of $\Delta$ equals the sign of $2a_a^D-1$, i.e.\ of
$a_a^D-\tfrac12$, whenever $P_T(D)>0$ (and $\Delta=0$ when $P_T(D)=0$, in which case
the decision is moot).

\emph{The identifiability equivalence.} ($\Leftarrow$) If the evidence determines the
event $\{a_a^D>\tfrac12\}$ --- for instance, a label-free reliability signal brackets
$a_a^D$ strictly on one side of $\tfrac12$ --- then by the identity it determines
$\sign\Delta$. ($\Rightarrow$) Conversely, if two target worlds are compatible with the
same evidence yet place $a_a^D$ on opposite sides of $\tfrac12$, the identity forces
their benefits to have opposite signs, which is precisely the unknowable configuration
of Theorem~\ref{thm:imp}; so identifiability of $\sign\Delta$ requires the evidence to
settle the above-chance question. \qedhere
\end{proof}

\begin{remark}[Ordinal, not cardinal --- and exactly one residual unknown]
\label{sec:t2-rem-ordinal}
The theorem changes the problem's type. Estimating either absolute risk $R_T(f_0)$ or
$R_T(f_a)$ without labels is the unsupervised-risk-estimation problem, known --- from
the conditional-independence analyses of Steinhardt and Liang and the divergence-based
bounds reviewed in the related-work discussion --- to be impossible without structural
assumptions.
The identity shows the adapt/freeze decision never needed those cardinal objects: it
needs one \emph{ordinal} bit --- is $f_a$ above chance on $D$ --- a strictly weaker
requirement. Everything else in the formula is label-free: $P_T(D)$ is an observable
mass, and the per-sample benefits off $D$ are identically zero. The residual burden is
exactly the bracketing of $a_a^D$ around $\tfrac12$, and the witness of
Theorem~\ref{thm:imp} shows that burden is real: there $f_0=\ind[x>0]$ and
$f_a=\ind[x<0]$ disagree almost everywhere, so $D$ is essentially the whole line, and
the two worlds set $a_a^D=0$ and $a_a^D=1$ respectively --- maximal disagreement about
the ordinal bit, with identical evidence. The witness is thus precisely a failure of
label-free bracketing on $D$, which is why the open problem of
Section~\ref{sec:t2-open} has the shape it has.
\end{remark}

The binary identity generalizes in both directions the deployment cares about ---
multiclass classification and regression --- and the generalizations are the results
that close the characterization half of the program.

\begin{theorem}[Multiclass $0/1$ loss]
\label{thm:disagree-mc}
For $Y\in\{1,\dots,K\}$ with $0/1$ loss, with $p_a=P_T(f_a{=}Y\mid D)$,
$p_0=P_T(f_0{=}Y\mid D)$, off $D$ $\delta=0$ and on $D$ $\E[\delta\mid D]=p_a-p_0$, so
$\Delta=P_T(D)(p_a-p_0)$ and $\sign\Delta=\sign(p_a-p_0)$.
\end{theorem}

\begin{proof}
As before $\delta=\ind[f_0\neq Y]-\ind[f_a\neq Y]$, and off $D$ the two indicators
coincide pointwise, so $\delta=0$ there. On $D$ the predictions differ, and --- this is
the multiclass novelty --- the conditional space splits into \emph{three} disjoint
events rather than two: $\{f_a=Y\}$ (the candidate is right, the incumbent necessarily
wrong since it predicts something different, $\delta=1-0=+1$), $\{f_0=Y\}$ (the
incumbent is right, the candidate wrong, $\delta=0-1=-1$), and $\{f_0\neq Y, f_a\neq
Y\}$ (both wrong, predicting two \emph{different} wrong classes, where
$\delta=1-1=0$). The two correctness events are still disjoint on $D$ --- the models
cannot both be right while disagreeing --- but they are no longer complementary:
$p_a+p_0\le1$ with strict inequality whenever both-wrong has mass. Taking the
conditional expectation over the three cells,
\[
\E[\delta\mid D]
=(+1)\,p_a+(-1)\,p_0+0\cdot\bigl(1-p_a-p_0\bigr)
=p_a-p_0 ,
\]
which one can equally read off the complement form
$\E[\delta\mid D]=(1-p_0)-(1-p_a)$. Total expectation over $\{D,D^c\}$ gives
$\Delta=P_T(D)\,\E[\delta\mid D]=P_T(D)(p_a-p_0)$, and nonnegativity of $P_T(D)$ gives
the sign identity. For $K=2$, exactly one predictor is correct on $D$, so
$p_0=1-p_a$, the both-wrong cell is empty, and the formula reduces to
$\Delta=P_T(D)(2a_a^D-1)$ --- Theorem~\ref{thm:disagree}. For $K\ge3$ both may err,
and the identity holds regardless, because the benefit counts only who is
\emph{correct}, not whether the other is wrong. \qedhere
\end{proof}

\begin{theorem}[Squared-error regression]
\label{thm:disagree-reg}
For $Y\in\R$ with squared loss, $\delta=(f_0-f_a)(f_0+f_a-2Y)$, so on $D$
$\E[\delta\mid D]=\E[(f_0-f_a)(f_0+f_a)\mid D]-2\,\E[(f_0-f_a)Y\mid D]$. The first term is
observable from unlabeled target; thus $\sign\Delta$ is label-free identifiable iff the
sign contribution of $\E[(f_0-f_a)Y\mid D]$ is certifiable---which holds under covariate
shift (importance-weighted source estimate of $\E[Y\mid X]$) and fails under concept
shift (re-entering Theorem~\ref{thm:imp}).
\end{theorem}

\begin{proof}
The pointwise identity is the difference of squares:
\[
\delta=(f_0-Y)^2-(f_a-Y)^2
=\bigl[(f_0-Y)-(f_a-Y)\bigr]\bigl[(f_0-Y)+(f_a-Y)\bigr]
=(f_0-f_a)\,(f_0+f_a-2Y).
\]
Off $D$ we have $f_0(x)=f_a(x)$, so the first factor vanishes and $\delta=0$; the
disagreement region again carries the whole benefit, and
$\Delta=P_T(D)\,\E[\delta\mid D]$ by total expectation. On $D$, expanding by linearity,
\[
\E[\delta\mid D]
=\E\bigl[(f_0-f_a)(f_0+f_a)\,\big|\,D\bigr]\;-\;2\,\E\bigl[(f_0-f_a)\,Y\,\big|\,D\bigr].
\]
The first term is a functional of the two \emph{known} maps $f_0,f_a$ and of
$\Law(X\mid D)$ only --- estimable from the unlabeled target stream to arbitrary
precision. The second term is the unique label-coupled object, so $\sign\Delta$ is
label-free identifiable exactly when that term's sign contribution can be certified
without target labels. Under covariate shift it can: writing $m(X)=\E[Y\mid X]$ (the
same function in source and target by the premise), the tower property and
$\sigma(X)$-measurability of $f_0-f_a$ give
\[
\E\bigl[(f_0-f_a)\,Y\,\big|\,D\bigr]
=\E\bigl[(f_0-f_a)\,\E[Y\mid X]\,\big|\,D\bigr]
=\E\bigl[(f_0-f_a)\,m(X)\,\big|\,D\bigr],
\]
where $m$ is estimable from \emph{labeled source} data (importance-weighted as in
Theorem~\ref{thm:pos}) and the remaining ingredients are target-observable. Under
concept shift the source regression function $m_S$ is the wrong $m$, the plug-in can
carry the wrong sign, and no label-free correction is available --- the configuration
re-enters Theorem~\ref{thm:imp}. \qedhere
\end{proof}

\subsection{Numerical Validation of the Identities}
\label{sec:t2-disagree-validation}

Algebraic identities invite the most direct kind of machine checking, and the
repository's validator
\texttt{experiments/kbound/theory\_validation/val\_thm5\_multiclass.py}
\cite{kbound2026repo} performs it at scale. For the multiclass identity it draws
$4000$ random trials --- class counts $K$ ranging over $\{3,\dots,11\}$, sample sizes
between $2000$ and $8000$, predictors of independently random accuracies so that
$\Delta$ takes both signs across trials --- and computes $\Delta$ two ways: directly
with labels (the mean per-sample benefit), and through the decomposition
$P_T(D)(p_a-p_0)$. As recorded in the paper, the identity holds to $10^{-16}$ with
$100\%$ sign agreement over the $4000$ trials --- and, decisively for the
beyond-binary claim, the both-wrong-on-$D$ event occurs in \emph{every} trial, so the
regime genuinely exercised is the one where $p_0\neq1-p_a$. The script additionally
verifies that the off-$D$ contribution to $\Delta$ is exactly zero, sample by sample.
For regression it checks the pointwise difference-of-squares identity to machine
precision over $3000$ trials (tolerance $10^{-9}$, with $100\%$ sign agreement between
$\Delta$ and $\E[\delta\mid D]$), and then runs the shift demonstration of
Theorem~\ref{thm:disagree-reg} end to end: with $40{,}000$ source and target points, a
$2.0$-standard-deviation covariate shift, and a source-only kernel-regression plug-in
for $m(X)$, the label-free certificate recovers the correct $\sign\Delta$ under
covariate shift and \emph{flips} sign under concept shift --- the predicted failure,
observed on cue. The validation does not, of course, prove the theorems (the proofs
above do); what it certifies is that the implemented quantities are the theorems'
quantities, with no sign conventions or conditioning slips between paper and code.

\subsection{Disagreement as the Decisive Evidence Coordinate}
\label{sec:t2-disagree-evidence}

The theorems of this section make a falsifiable prediction about \emph{evidence
design}: if the entire benefit is carried by $D$, then among the components of the
evidence vector $Z$, the ones measuring disagreement between the frozen and candidate
models should be the most decision-relevant --- not marginally, but structurally. The
ablation study of Chapter~\ref{ch:exp1} tests exactly this by deleting evidence
features one at a time and re-measuring the wrapper's regret-to-oracle: removing the
detector-disagreement feature raises regret from $0.037$ to $0.094$, a $2.6\times$
degradation and by far the largest of any feature (the next-largest deletions move
regret to $0.043$ and $0.039$). The paper draws the connection explicitly: detector
disagreement is the single most important evidence feature in the ablation, directly
supporting that the disagreement region carries the decision. The theory says where
the information lives; the ablation confirms the estimator starves without it.

\section{The Open Problem: Label-Free Bracketing}
\label{sec:t2-open}

Everything proved so far reduces the label-free adapt/freeze decision to a single
remaining object: an ordinal comparison on the observable region $D$. The paper is
exact about what is closed and what is not, and states the open piece as a conjecture.
We reproduce it verbatim, and then spend an honest page on why it is genuinely open.

\begin{conjecture}[Label-free bracketing; \textbf{the remaining open piece}]
\label{conj:gen}
Theorems~\ref{thm:disagree-mc}--\ref{thm:disagree-reg} \emph{close the characterization}:
$\sign\Delta$ equals an ordinal accuracy comparison on the observable region $D$. What
remains open is the label-free \emph{bracketing} of that comparison ($p_a$ vs $p_0$ for
$K\ge3$, where $p_0\neq1-p_a$ makes it a two-sided test; or the sign of
$\E[(f_0-f_a)Y\mid D]$ for regression) without target labels---a reliability-model
assumption, as in the binary case.
\end{conjecture}

\subsection{Why Current Techniques Do Not Resolve It}
\label{sec:t2-open-why}

To see the precise difficulty, fix the multiclass case and ask for the object the
conjecture demands: a label-free functional of the observables --- the unlabeled target
stream on $D$, the two prediction maps, and whatever source-side calibration data one
allows --- that outputs an interval $[L,U]\ni p_a-p_0$ with $1-\alpha$ coverage, such
that committing when $0\notin[L,U]$ is valid \emph{without} a premise that is itself
unverifiable from the same observables. Every known route founders on one of two
obstructions.

\emph{The information obstruction.} Unconditionally --- over all label laws on $D$
consistent with the observed unlabeled data --- the task is hopeless, and this is a
theorem, not a suspicion: the witness of Theorem~\ref{thm:imp}, restricted to $D$,
exhibits two worlds with identical evidence laws and $(p_a,p_0)$ equal to $(1,0)$ and
$(0,1)$ respectively. Any bracket with valid coverage in both worlds must contain both
$+1$ and $-1$, hence $0$, hence never commits. This is the two-point method of Le~Cam
in its simplest form \cite{lecam1973convergence,tsybakov2009introduction}: the
unlabeled likelihood does not separate the hypotheses, so additional structure is
\emph{necessary}, and the conjecture is properly read as a question about the
\emph{weakest} usable structure rather than about removing structure altogether. The
binary case already had this character --- there the residual assumption is a
reliability signal bracketing $a_a^D$ against $\tfrac12$ --- but the multiclass case is
strictly harder in a concrete sense the paper flags in the conjecture itself: since
$p_0\neq1-p_a$, one scalar no longer determines the comparison, and the test becomes
two-sided in two unknown conditional accuracies linked only by the inequality
$p_a+p_0\le1$.

\emph{The technique obstruction.} Each available tool discharges the bracketing only by
importing a premise of the same epistemic type as the one being discharged.
Importance weighting (Theorem~\ref{thm:pos}) certifies the regression term
$\E[(f_0-f_a)Y\mid D]$ --- but only under covariate shift, which fixes the unobservable
conditional $P(Y\mid X)$ by fiat and is untestable from unlabeled data. Confidence-based
bracketing (the route formalized below) bounds $p_h$ by the mean calibrated confidence
$q_h$ --- but only under a calibration-transfer premise, a bound $\eta$ on the
\emph{target-side} miscalibration on $D$, which is again a quantity no label-free
procedure can estimate; assuming $\eta$ small is assuming a reliability model.
Label-shift machinery \cite{lipton2018detecting} identifies shifted class priors under
stable class-conditionals and an invertible confusion matrix --- structural premises
once more, and even granted them, priors do not by themselves yield the
\emph{conditional-on-$D$} accuracies the comparison needs. Empirical regularities such
as agreement-based accuracy extrapolation are calibrated correlations, not guarantees,
and nothing prevents the adversarial pairs of Theorem~\ref{thm:imp} from sitting
exactly where the regularity breaks. In every case the assumption is doing the work
precisely on the coordinate the data cannot see; the open question is whether some
premise exists that is weaker, or --- more interestingly --- \emph{self-certifying}, in
the sense that its failure is itself label-free detectable, so that the bracket could
widen automatically where the premise degrades.

\subsection{What Is Actually Proved}
\label{sec:t2-open-partial}

The honest ledger has four entries on the proved side. First, the
\emph{characterization} is closed unconditionally: Theorems~\ref{thm:disagree}
through~\ref{thm:disagree-reg} are exact identities requiring no assumptions beyond
the loss structure, and they localize the entire decision to $D$. Second, the binary
case needs only a \emph{one-sided} bracket of a single scalar against the fixed
threshold $\tfrac12$ --- a structurally easier object than the multiclass two-sided
comparison, which is why the binary instantiation is the one deployed throughout the
experiments. Third, the regression term is certified under covariate shift
(Theorem~\ref{thm:disagree-reg} with Theorem~\ref{thm:pos}). Fourth, the paper proves
a \emph{conditional} closure of the conjecture under a quantitative reliability model
--- bounded calibration transfer --- with an exact and provably tight threshold. We
restate it faithfully.

\begin{proposition}[Label-free bracketing under bounded calibration transfer]
\label{sec:t2-prop-caltransfer}
On the disagreement region $D$, let $\mathrm{conf}_h(x)\in[0,1]$ be a source-calibrated
confidence map for predictor $h\in\{f_0,f_a\}$, and suppose its target miscalibration
on $D$ is bounded by $\eta$:
$|\Pr_T(h(X){=}Y\mid \mathrm{conf}_h{=}c,\,X\in D)-c|\le\eta$ for all $c$
(\emph{calibration transfer}). Let $q_h:=\E_T[\mathrm{conf}_h(X)\mid D]$ be the
label-free mean confidence. Then $|q_h-p_h|\le\eta$, hence
$p_a-p_0\in[(q_a-q_0)-2\eta,\,(q_a-q_0)+2\eta]$, and
\[
|q_a-q_0|>2\eta \;\Longrightarrow\; \sign\Delta=\sign(q_a-q_0),
\]
so the multiclass/regression benefit sign is identified label-free; when
$|q_a-q_0|\le2\eta$ the bracket contains $0$ and the rule must \textsc{abstain}.
\end{proposition}

\begin{proof}
By the tower property,
$p_h=\E_T\bigl[\Pr_T(h{=}Y\mid \mathrm{conf}_h,D)\mid D\bigr]$; the
calibration-transfer bound says the inner conditional probability is within $\eta$ of
$\mathrm{conf}_h$ pointwise, so integrating preserves the bound and $|p_h-q_h|\le\eta$
for each $h$. Differencing the two $\eta$-balls gives
$|(p_a-p_0)-(q_a-q_0)|\le2\eta$, the stated bracket. If $|q_a-q_0|>2\eta$, the bracket
lies strictly on one side of $0$, so $\sign(p_a-p_0)=\sign(q_a-q_0)$; by
Theorem~\ref{thm:disagree-mc} this equals $\sign\Delta$. If $|q_a-q_0|\le2\eta$ the
bracket straddles $0$ and certifies nothing, so the certificate logic of
Theorem~\ref{thm:cert} mandates abstention. \qedhere
\end{proof}

The threshold $2\eta$ is \emph{exact}, as the paper notes: when $|q_a-q_0|\le2\eta$ an
adversary inside the two $\eta$-balls can place $p_a,p_0$ on either side of equality,
re-entering the unknowable regime of Theorem~\ref{thm:imp} --- so the proposition's
abstention region cannot be shrunk under its own premise. The validator
(\texttt{val\_conj1\_caltransfer.py}, results in
\texttt{results\_conj1\_caltransfer.json} \cite{kbound2026repo}) sweeps both $\eta$ and
the true gap: with a true gap of $0.2$ ($p_a=0.6$, $p_0=0.4$), every committed decision
is sign-correct ($100\%$ recovery among committed, $0$ wrong commitments) for all
$\eta\in\{0,0.02,0.05,0.08\}$ where the gap exceeds $2\eta$; at the knife edge
$\eta=0.1$ (gap exactly $2\eta$) the finite-sample rule still commits on $52.45\%$ of
trials --- all of them sign-correct, with the wrong-commit rate still $0$ --- and for
$\eta\ge0.15$ commitment falls to $0$, full abstention. The dual gap sweep at
$\eta=0.05$ shows the same wall from the other side (no commitments below gap $0.10$,
full commitment from gap $0.12$ up), and the bracket
$[(q_a-q_0)-2\eta,(q_a-q_0)+2\eta]$ contains the true gap on $100\%$ of trials. The
headline fields record minimum sign recovery $1.0$ and maximum wrong-commit rate $0.0$
whenever the gap exceeds $2\eta$ --- the guarantee, observed.

\subsection{What a Resolution Would Buy}
\label{sec:t2-open-payoff}

We state plainly that the conjecture is \emph{not} solved here, and that
Proposition~\ref{sec:t2-prop-caltransfer} does not solve it: it trades the open
assumption for a quantified one, which is progress in precision, not in kind. A
genuine resolution would change the system's standing in either direction. A
\emph{positive} resolution --- a bracketing scheme whose premise is label-free
verifiable, or a proof that some natural observable (margin spectra on $D$, agreement
geometry among auxiliary predictors, weight diagnostics) suffices --- would make the
entire trichotomy assumption-free given evidence: \KGA's residual reliability-model
assumption would vanish, the knowable region of the phase diagram would extend to
every shift with nonvacuous brackets, and evidence design would reduce to estimating
two scalars on an observable set. A \emph{negative} resolution --- a lower bound
showing every label-free bracket valid over a natural shift class must contain $0$
unless a premise of calibration-transfer strength holds --- would be just as valuable:
it would certify that the conditional results of this chapter are the strongest
possible, closing the program with matching upper and lower bounds the way
Theorem~\ref{thm:imp} closed the unconditional question. Either outcome turns the last
assumption in the stack into a theorem. Until then, the honest summary is the paper's
own: the characterization is closed; the label-free estimability of the ordinal
comparison is the standing assumption; the unconditional bracketing is the one residual
open piece.

\section{Synthesis: The Certified Map of Regimes}
\label{sec:t2-synthesis}

Chapters~\ref{ch:theory1} and~\ref{ch:theory2} together draw one map. The
impossibility theorem fences off the region where no label-free rule may commit; the
certificate gives the safe-passage protocol everywhere else; the covariate-shift and
disagreement-region results chart where the passage provably opens; and the conjecture
marks the one stretch of coastline still unsurveyed. Table~\ref{sec:t2-tab-map}
assembles the chapter's results with their guarantees and the machine-checked evidence
behind each, with the experimental chapters and Appendix~\ref{app:tables} holding the
full per-condition records.

\begin{table}[t]\centering\small
\begin{tabular}{p{0.20\textwidth}p{0.15\textwidth}p{0.27\textwidth}p{0.28\textwidth}}
\toprule
Regime & Result & Guarantee & Validator / evidence \\
\midrule
Matched evidence (unknowable) & Thm~\ref{thm:imp}; Cor.~\ref{sec:t2-cor-abstain} &
no label-free rule certifies both worlds; any valid certificate abstains with
probability $\ge1-2\alpha$ &
clean witness: $100\%$ abstention, forced-commit regret $0.51$
(Ch.~\ref{ch:exp1}) \\
\addlinespace
Risk-aligned evidence, batch & Thm~\ref{thm:cert}; Cor.~\ref{sec:t2-cor-n} &
false-adapt $\le\alpha$ and false-freeze $\le\alpha$; commits once
$n\gtrsim\sigma^2\Delta^{-2}\log(1/\alpha)$ &
mixed regime: false-adapt $0.065$ at $\alpha=0.1$; stress grid: $0\%$ false-adapt
(Chs.~\ref{ch:exp1}--\ref{ch:exp2}) \\
\addlinespace
Risk-aligned evidence, streaming & Thm~\ref{thm:evalue} &
anytime-valid: error $\le\alpha$ simultaneously over all stopping times &
worst anytime false-adapt $0.0626\le0.1$
(\texttt{val\_thm3\_evalue\_results.json}) \\
\addlinespace
Covariate shift & Thm~\ref{thm:pos} &
$\Delta$ and $\sign\Delta$ identifiable from labeled source $+$ unlabeled target &
regression demo: sign correct under covariate shift, flips under concept shift
(\texttt{val\_thm5\_multiclass.py}) \\
\addlinespace
Disagreement region & Thms~\ref{thm:disagree}--\ref{thm:disagree-reg} &
$\sign\Delta$ $=$ ordinal accuracy comparison on observable $D$ (exact identity) &
identity to $10^{-16}$, $100\%$ sign agreement over $4000$ trials
(\texttt{val\_thm5\_multiclass.py}) \\
\addlinespace
Multiclass / regression bracketing & Conj.~\ref{conj:gen} (open);
Prop.~\ref{sec:t2-prop-caltransfer} (conditional) &
sign certified iff $|q_a-q_0|>2\eta$ under calibration transfer; unconditional case
open &
$100\%$ sign recovery, $0$ wrong commitments when gap $>2\eta$
(\texttt{results\_conj1\_caltransfer.json}) \\
\bottomrule
\end{tabular}
\caption{The certified map of regimes assembled by Chapters~\ref{ch:theory1}
and~\ref{ch:theory2}: each row pairs a regime with its governing result, the exact form
of its guarantee, and the machine-checked validator or experiment that exercises it.}
\label{sec:t2-tab-map}
\end{table}

Three through-lines deserve a final statement. First, every positive result in this
chapter is \emph{conditional in a way the theory itself prices}. The certificate is
valid given a valid radius; the radius is valid given exchangeability of the
calibration bank; covariate-shift identifiability is valid given an untestable premise
on $P(Y\mid X)$; the bracketing is valid given calibration transfer. None of this is a
defect to be hidden --- it is the content of Theorem~\ref{thm:imp} propagating through
the constructions, and the certificate's distinctive virtue is that where its premises
fail detectably (heavy importance weights, wide residual quantiles, sub-threshold
confidence gaps) the failure surfaces as a \emph{wider radius}, and the rule degrades
to abstention rather than to silent error. Second, the disagreement region is the
load-bearing structure connecting theory to practice: it is where the benefit lives
(Theorems~\ref{thm:disagree}--\ref{thm:disagree-reg}), why the empirical-Bernstein
radius is tight (Remark~\ref{sec:t2-rem-variance}), and which evidence coordinate the
deployed estimator cannot do without (the $2.6\times$ ablation of
Chapter~\ref{ch:exp1}). Third, the guarantees proved here are exactly the ones the
deployed system claims --- no more. Chapter~\ref{ch:kga} assembles
$Z\to\widehat\Delta\to\varepsilon\to\{\textsc{adapt},\textsc{freeze},\textsc{abstain}\}$
into the \KGA\ wrapper with the assumptions of this chapter stated as the algorithm's
assumptions, and Chapters~\ref{ch:exp1} and~\ref{ch:exp2} then measure the certificate
against worlds engineered to reward it, punish it, and --- on the constructed witness
--- defeat it, which is the only honest order of business after a chapter of
conditional guarantees.
